\chapter{点云几何特征检测实验结果及分析}
\label{chap:experiment}

本章旨在通过系统性的实验验证本文所提出的 NECDNet 在三维点云边缘与角点检测任务中的有效性与鲁棒性。实验主要基于本文构建的 ABC-Corner 数据集展开。首先，本章将详细阐述 ABC-Corner 数据集的构建过程、实验的硬件环境、参数配置以及针对数据不平衡问题所采用的评价指标与训练策略。随后，将 NECDNet 与当前主流的传统几何分析方法及深度学习方法进行多维度的定性与定量对比，并重点分析法向量估计精度对检测结果的显著影响。此外，为了验证网络结构设计的合理性，本章还对特征维度、多尺度参数选择及数据增强策略进行了详尽的消融实验分析。最后，本章对实验结果进行了总结与讨论。

\section{实验环境及配置}

\subsection{软硬件环境与参数设置}
\label{sec:experiment setup}

为了确保实验结果的可复现性，本节首先明确实验的软硬件环境。本次实验的计算平台采用一台配备 NVIDIA RTX 3090 图形处理单元（显存 256 GB）的高性能工作站，运行环境为 Windows 10 操作系统。算法的开发与实现基于 PyTorch 深度学习框架~\cite{paszke_pytorch_2019}，集成开发环境为 PyCharm。

在网络模型的训练阶段，超参数的设定对于模型的收敛速度和最终性能至关重要。为了平衡计算效率与显存占用，将批次大小（Batch Size）设定为 20,000，共进行 64 个 Epoch 的迭代训练。为了加速大规模点云数据的读取与预处理，数据加载器的并行工作线程数被配置为 8。模型优化采用 Adam 优化器，该优化器结合了动量法和自适应学习率的优点，能够有效应对稀疏梯度问题。其动量参数分别设定为 $\beta_1 = 0.9$ 和 $\beta_2 = 0.999$，初始学习率设置为 0.001，配合余弦退火策略进行学习率衰减，以确保网络在训练初期的快速收敛以及后期的稳定优化。

针对三维点云数据中普遍存在的几何特征分布不平衡问题，即面点数量远多于边点，而角点最为稀缺，传统的交叉熵损失函数往往会导致模型被大量简单样本主导。为此，本文采用了 Focal Loss~\cite{lin_focal_2020} 作为损失函数，并将聚焦参数 $\gamma$ 设为 2.5。该参数能够动态降低易分类样本的权重，使得模型在训练过程中更加关注那些难以分类的边缘与角点样本，从而有效缓解类别不平衡带来的负面影响。

\subsection{评价指标设计}

为了全面且客观地评估模型的性能，本文选取精确率（Precision）和召回率（Recall）作为主要评价指标，同时辅以准确率（Accuracy）。此外，为了与现有的 SOTA 方法进行公平对比，还计算了马修斯相关系数（MCC）、F1 分数以及交并比（IoU）。

对于纯边缘检测任务，直接采用标准的二分类指标公式。而对于边缘与角点的联合检测任务，这是一个典型的多分类问题，且存在严重的类别不平衡。若直接使用全局准确率，模型倾向于将所有点预测为面点即可获得极高的分数，这显然是不合理的。为此，本文采用宏平均（Macro-averaging）策略来计算 Precision、Recall 和 F1 分数。具体而言，首先基于各类别的真正例（TP）、假正例（FP）、真负例（TN）和假负例（FN）统计量分别计算该类别的指标，然后对所有类别的指标取算术平均。这种计算方式赋予了每个类别同等的权重，确保了稀疏类别（如角点）的检测效果能够被充分反映，避免了评价结果被数量占优的面点所主导。

\section{对比实验结果及分析}

为了全面评估 NECDNet 的性能，本节将其与现有的 SOTA 方法进行了对比。对比实验涵盖两个层面：一是针对边缘与角点联合检测的多分类任务，对比方法包括基于协方差分析的 CA 方法~\cite{bazazian_fast_2015}、两阶段检测法 TECD~\cite{ahmed_edge_2018} 以及深度学习方法 PIE-Net~\cite{wang_pienet_2020a}；二是针对纯边缘检测的二分类任务，对比方法包括 PCEDNet~\cite{himeur_pcednet_2022} 和 BoundED~\cite{bode_bounded_2023}。

\subsection{边缘与角点联合检测结果分析}

在边缘与角点的联合检测任务中，表~\ref{tab:corner_result} 展示了各方法在 ABC-Corner 数据集上的定量评估结果。数据表明，NECDNet 在 Precision、Recall、F1 分数、Accuracy、IoU 及 MCC 等所有关键指标上均显著优于对比方法。

具体分析各方法的表现，传统几何分析方法 CA 和 TECD 的各项指标相对较低。CA 方法主要依赖单一尺度的几何统计量（如协方差矩阵特征值）进行阈值判定。这种硬阈值策略难以适应不同物体间巨大的几何差异，例如在处理薄壁物体时，单一尺度的邻域分析容易跨越物体的两个表面，导致大量的误检。TECD 方法虽然采用了两步策略，即先检测边缘再基于边缘检测角点，但其角点检测的准确性高度依赖于第一阶段边缘提取的质量，且存在复杂的阈值参数依赖，导致其泛化能力不足，在复杂拓扑结构下的鲁棒性较差。

相比之下，基于深度学习的 PIE-Net 虽然在性能上优于传统方法，但其 F1 分数仍低于本文方法。PIE-Net 采用了两个独立的网络分支分别处理边缘和角点，缺乏对两者特征相关性的联合建模。这种设计割裂了边缘与角点在几何上的内在联系，导致在某些模糊区域容易产生混淆，且存在较高的假阳性率。本文提出的 NECDNet 通过端到端的特征学习和多尺度特征融合，能够更有效地捕捉局部几何结构。特别是当使用从网格模型（Mesh）估计的精确法向量作为输入时，模型的 F1 分数达到 0.815，MCC 达到 0.943，展现了极高的分类一致性。

\begin{table*}[ht]
    \centering
    \small
    \setlength{\tabcolsep}{17.6pt}
    \begin{threeparttable}
    \caption{ABC-Corner 数据集上各边缘与角点检测方法的宏平均（Macro-averaging）指标评分。}
    \label{tab:corner_result}
    \begin{tabular}{lcccccc}
        \toprule
        Method & Precision($\uparrow$) & Recall($\uparrow$) & F1($\uparrow$) & Accuracy($\uparrow$) & IoU($\uparrow$) & MCC($\uparrow$) \\
        \midrule
        CA~\cite{bazazian_fast_2015} & 0.415 & 0.616 & 0.420 & 0.616 & 0.346 & 0.276 \\
        TECD~\cite{ahmed_edge_2018} & 0.472 & 0.642 & 0.513 & 0.642 & 0.438 & 0.526 \\
        PIE-NET~\cite{wang_pienet_2020a} & 0.615 & 0.814 & 0.686 & 0.935 & 0.577 & 0.718 \\
        NECDNet (Ours)  & 0.646 & 0.857 & 0.715 & 0.983 & 0.616 & 0.833 \\
        NECDNet$\dagger$ (Ours)  & \textbf{0.765} & \textbf{0.933} & \textbf{0.815} & \textbf{0.995} & \textbf{0.743} & \textbf{0.943} \\
        \bottomrule
    \end{tabular}
    \begin{tablenotes}
        \footnotesize
        \item PIE-Net 的数据基于其作者提供的评估结果整理得出。
        
        \item $\dagger$ 表示使用从 Mesh 模型估计的法向量。 
     \end{tablenotes}
    
    \end{threeparttable}
    
\end{table*}

定性分析进一步印证了上述结论。如图~\ref{fig:edge_corner_result} 所示的可视化结果显示，NECDNet 的检测结果与真值（Ground Truth）最为接近，能够清晰地勾勒出物体的骨架结构并准确定位角点。CA 和 TECD 方法的鲁棒性较差，出现了较多的噪点；PIE-Net 虽然检测出了大部分特征，但在角点定位上不够精准。PCEDNet 虽然在边缘检测上表现良好，但由于无法检测角点，在图示中缺失了红色角点标记。为了更细致地探究各方法在不同几何拓扑下的表现，表~\ref{tab:metrics_result_of_four_obj} 列出了针对图~\ref{fig:edge_corner_result} 中四个典型复杂物体的分类别详细指标。结果显示，NECDNet 在最具挑战性的“角点”类别上取得了压倒性优势，其角点 F1 分数大幅领先于第二名，证明了本文提出的多尺度特征融合模块在捕捉稀疏拓扑特征方面的有效性。

\begin{figure*}[t]
      \centering
       \includegraphics[width=\linewidth]{NECDNet/result.jpg}
       \caption{ABC-Corner 数据集上各方法的边缘与角点检测结果可视化对比。本文方法（Ours）取得了最佳的检测效果。CA 和 TECD 缺乏鲁棒性，PIE-Net 检测精度不足，PCEDNet 表现出良好的边缘检测能力但无法检测角点。}
       \label{fig:edge_corner_result}
\end{figure*}

\begin{table*}[ht]
        % \centering
        \small
        \setlength{\tabcolsep}{2.6pt}
        \caption{图~\ref{fig:edge_corner_result} 中四个点云对象的分类别详细指标评分。}
        \label{tab:metrics_result_of_four_obj}
        \begin{threeparttable}
            \begin{tabular}{lccc ccc ccc ccc ccc c}
            \toprule
            Method & \multicolumn{3}{c}{Precision($\uparrow$)} & \multicolumn{3}{c}{Recall($\uparrow$)} & \multicolumn{3}{c}{F1($\uparrow$)} & \multicolumn{3}{c}{Accuracy($\uparrow$)} & \multicolumn{3}{c}{IoU($\uparrow$)} & MCC($\uparrow$) \\
            \cmidrule(lr){2-4} \cmidrule(lr){5-7} \cmidrule(lr){8-10} \cmidrule(lr){11-13} \cmidrule(lr){14-16} 
             & Face & Edge & Corner & Face & Edge & Corner & Face & Edge & Corner & Face & Edge & Corner & Face & Edge & Corner &  \\
            \midrule
            CA~\cite{bazazian_fast_2015} & 0.975 & 0.303 & 0.001 & 0.582 & 0.614 & 0.303 & 0.729 & 0.405 & 0.002 & 0.582 & 0.614 & 0.303 & 0.573 & 0.254 & 0.001 & 0.184 \\
            TECD~\cite{ahmed_edge_2018} & 0.986 & 0.705 & 0.062 & 0.980 & 0.762 & 0.158 & 0.983 & 0.732 & 0.089 & 0.980 & 0.762 & 0.158 & 0.967 & 0.578 & 0.046 & 0.707 \\
            PIE-NET~\cite{wang_pienet_2020a} & 0.995 & 0.673 & 0.278 & 0.950 & 0.898 & 0.964 & 0.972 & 0.770 & 0.432 & 0.950 & 0.898 & 0.964 & 0.946 & 0.626 & 0.275 & 0.759 \\
            NECDNet (Ours) & \textbf{0.999}& 0.970 & 0.519 & 0.998 & \textbf{0.987}& 0.921 & \textbf{0.999}& 0.978 & 0.664 & 0.998 & \textbf{0.987}& 0.921 & 0.997 & 0.957 & 0.496 & 0.972 \\
            NECDNet$\dagger$ (Ours) & \textbf{0.999}& \textbf{0.994}& \textbf{0.540}& \textbf{1.000}& 0.974 & \textbf{0.987}& \textbf{0.999}& \textbf{0.984}& \textbf{0.698}& \textbf{1.000}& 0.974 & \textbf{0.987}& \textbf{0.999}& \textbf{0.968}& \textbf{0.536}& \textbf{0.983}\\
            \bottomrule
        \end{tabular}
        \begin{tablenotes}
            \footnotesize
            \item PIE-Net 的数据基于其作者提供的评估结果整理得出。
            
            \item $\dagger$ 表示使用从 Mesh 模型估计的法向量。 
         \end{tablenotes}
        \end{threeparttable}
        
    \end{table*}

\subsection{纯边缘检测结果分析}

为了验证模型在通用边缘检测任务上的有效性，表~\ref{tab:edge_result} 展示了在 ABC 数据集上仅关注边缘二分类任务时的对比结果。由于 PCEDNet 和 BoundED 未开源代码，我们引用了其论文中的数据进行对比。

实验结果表明，NECDNet 在该任务上同样表现出色。在使用 Mesh 法向量的情况下（NECDNet-2c$\dagger$），其 F1 分数高达 0.956，显著优于 PCEDNet (0.822) 和 BoundED (0.875)。即便是使用基于局部拟合平面的估计法向量（NECDNet-2c），其性能也具备很强的竞争力，尤其是在召回率上表现优异，说明模型能够尽可能多地检出潜在边缘。这表明 NECDNet 不仅是一个优秀的特征点检测器，也是一个鲁棒的边缘提取器，能够适应不同的应用需求。

\begin{table*}[ht]
    \centering
    \small
    \begin{threeparttable}
    \setlength{\tabcolsep}{17pt}
    \caption{ABC 数据集上各边缘检测方法的指标评分。}
    \label{tab:edge_result}
    \begin{tabular}{lcccccc}
        \toprule
        Method & Precision($\uparrow$) & Recall($\uparrow$) & F1($\uparrow$) & Accuracy($\uparrow$) & IoU($\uparrow$) & MCC($\uparrow$) \\
        \midrule
        PCEDNet-2c$\dagger$~\cite{himeur_pcednet_2022} & 0.735 & 0.984 & 0.822 & 0.970 & 0.597 & 0.808 \\
        BoundED-2c~\cite{bode_bounded_2023} & \textbf{0.959} & 0.857 & 0.875 & 0.987 & 0.778 & 0.869 \\
        NECDNet-2c (Ours) & 0.708 & 0.988 & 0.825 & 0.980 & 0.702 & 0.827 \\
        NECDNet-2c$\dagger$ (Ours) & 0.922 & \textbf{0.992} & \textbf{0.956} & \textbf{0.996} & \textbf{0.915} & \textbf{0.954} \\
        \bottomrule
    \end{tabular}
    \begin{tablenotes}
        \footnotesize
         \item PCEDNet 和 BoundED 的数据分别整理自文献~\cite{himeur_pcednet_2022} 和~\cite{bode_bounded_2023}。
        
        \item $\dagger$ 表示使用从 Mesh 模型估计的法向量。 
     \end{tablenotes}
     \end{threeparttable}
    
\end{table*}

\subsection{讨论：法向量估计精度对检测的影响}

在深入分析实验结果后，本文发现法向量的精度是影响检测性能的关键因素。NECDNet 的核心输入特征——包括法向角差、法向分布及投影法向模长——均直接依赖于法向量的计算。因此，精确的法向量估计对于提升 NECDNet 的性能至关重要。

为了量化这一影响，本文分别测试了基于点云局部表面最小二乘拟合（Least Squares Fitting）和基于网格模型（Mesh）几何信息两种法向量估计方式。从几何原理上讲，最小二乘拟合假设邻域点位于一个光滑平面上。然而，在物体的尖锐边缘和角点附近，这一假设被打破：邻域点横跨了两个或多个法向截然不同的平面。此时，拟合算法往往会产生平滑过渡的“平均化”误差法向，而非真实的表面法向。

如图~\ref{fig:limitation} 的对比实验所示，这种方向误差会导致显著的分类混淆。在图~\ref{fig:limitation}(c) 中，由于法向的平滑，原本尖锐的边缘特征在特征空间中变得模糊，导致网络将部分边缘点误判为角点，或者将角点误判为边缘点。相反，利用 Mesh 计算出的法向量能够精确反映物体的几何拓扑（图~\ref{fig:limitation}(b)），帮助网络准确捕捉微小的曲率变化，从而实现近乎完美的分类效果。这一发现表明，虽然 NECDNet 具有强大的特征提取能力，但在实际应用中，高质量的法向信息预处理是发挥其性能上限的必要条件。

\begin{figure*}[htbp]
    \centering
    \includegraphics[width=\linewidth]{NECDNet/limitation.jpg}
    \caption{法向量估计精度对检测结果的影响分析：(a) 检测结果真值（Ground Truth）；(b) 使用从 Mesh 估计的法向量得到的检测结果；(c) 使用从点云估计的法向量得到的检测结果。可以看出，使用点云估计的法向量会导致更多的边缘点被误分类为角点。}
    \label{fig:limitation}
\end{figure*}

\section{消融实验分析}
\label{sec:ablation_study}

为了验证 NECDNet 各个模块设计的合理性，本节基于 ABC-Corner 数据集开展了多维度的消融实验。实验设定随机抽取 100 个模型用于训练，另外 100 个用于测试，控制变量以确保结论的可靠性。

\subsection{特征维度的影响}

特征描述子的选择直接决定了网络输入信息的丰富程度与有效性。在构建特征向量时，除了最终选用的 6 个标量特征（$\kappa_i, \bar{d_i}, r_i, \mu_{\theta,i}, \sigma^2_{\theta,i}, \bar{\text{proj}_i}$）外，实验还尝试引入了 3 维 PCA 特征值向量 $\mathbf{\sigma}_i$ 和 3 维法向量差均值向量 $\bar{\Delta \mathbf{n}_i}$。表~\ref{tab:ablation_study} 中的“特征维度”部分展示了不同组合的性能差异。

实验结果显示，当模型仅使用 6 个标量特征时，各项评价指标反而达到了最高值。引入向量特征后，性能未升反降。分析其原因，$\mathbf{\sigma}_i$ 和 $\bar{\Delta \mathbf{n}_i}$ 虽然包含方向信息，但在无序点云网络中，这种绝对方向信息可能引入了对旋转的敏感性或数据噪声。由于点云网络（如 PointNet++ 结构）通常具有旋转不变性或对局部坐标系的依赖，直接输入全局或不稳定的方向向量可能增加了输入数据的维度和复杂性，导致“维度灾难”效应，反而干扰了网络对核心几何不变量的学习。因此，NECDNet 最终选择更加紧凑且鲁棒的 6 维标量特征组合。

\subsection{多尺度参数的选择}

多尺度特征融合是解决点云密度不均和特征尺度变化问题的关键。本节对尺度数量（Number of scales）和尺度范围（邻域大小 $k$）进行了探索。实验设定 $k$ 的取值范围覆盖 8 到 128，尺度数量从 2 增加至 8。表~\ref{tab:ablation_study} 中的“尺度选择”部分表明，随着尺度数量的增加，模型性能呈现先上升后下降的趋势。

当选择 4 个尺度（$k=8, 16, 32, 64$）时，模型取得了最佳的 F1 分数和 IoU。当尺度数量过少（如 2 个）时，网络的感受野受限，难以提取宏观的几何结构信息；而当尺度数量过多（如 6 个或 8 个）时，不仅计算开销显著增加，而且大尺度的邻域（如 $k=128$）引入了过多的无关背景噪声，平滑了局部的尖锐特征，从而抑制了网络的判别能力。因此，4 个尺度的设置在特征丰富度与噪声控制之间取得了最佳平衡。

\begin{table*}[ht]
    \centering
    \small
    \setlength{\tabcolsep}{7pt}
    \caption{关于特征维度选择和多尺度参数选择的消融实验结果。}
    \label{tab:ablation_study}
    \begin{threeparttable}
    \begin{tabular}{llcccccc}
            \toprule
            &  & Precision($\uparrow$) & Recall($\uparrow$)  & F1($\uparrow$) & Accuracy($\uparrow$) & IoU($\uparrow$) & MCC($\uparrow$) \\
            \midrule
   
            \multirow{6}{*}{特征维度}
            & $(\hat{\mathbf{\sigma}_i}, \kappa_i, \bar{d_i}, r_i)$ & 0.647 & 0.847 & 0.654 & 0.977 & 0.601 & 0.790 \\
            & $(\bar{\Delta \mathbf{n}_i}, \mu_{\theta,i}, \sigma^2_{\theta,i}, \bar{\text{proj}_i})$ & 0.680 & 0.922 & 0.719 & 0.986 & 0.656 & 0.886 \\
            & $(\kappa_i, \bar{d_i}, r_i, \mu_{\theta,i}, \sigma^2_{\theta,i}, \bar{\text{proj}_i})$ & \textbf{0.754} & \textbf{0.946} & \textbf{0.805} & \textbf{0.996} & \textbf{0.735} & \textbf{0.946} \\
            & $(\kappa_i, \bar{d_i}, r_i, \bar{\Delta \mathbf{n}_i}, \mu_{\theta,i}, \sigma^2_{\theta,i}, \bar{\text{proj}_i})$ & 0.714 & 0.935 & 0.758 & 0.994 & 0.693 & 0.929 \\
            & $(\hat{\mathbf{\sigma}_i}, \kappa_i, \bar{d_i}, r_i, \mu_{\theta,i}, \sigma^2_{\theta,i}, \bar{\text{proj}_i})$ & 0.739 & 0.946 & 0.793 & 0.995 & 0.720 & 0.934 \\
            & $(\hat{\mathbf{\sigma}_i}, \kappa_i, \bar{d_i}, r_i, \bar{\Delta \mathbf{n}_i}, \mu_{\theta,i}, \sigma^2_{\theta,i}, \bar{\text{proj}_i})$ & 0.743 & 0.940 & 0.797 & 0.994 & 0.722 & 0.929 \\
           
            \midrule
           
            \multirow{4}{*}{尺度选择}
            & (8, 16) & 0.704 & 0.913 & 0.751 & 0.993 & 0.680 & 0.910 \\
            & (8, 16, 32, 64) & \textbf{0.743} & \textbf{0.940} & \textbf{0.797} & 0.994 & \textbf{0.722} & 0.929 \\
            & (8, 16, 32, 48, 80, 128) & 0.723 & 0.927 & 0.770 & \textbf{0.995} & 0.702 & \textbf{0.932} \\
            & (8, 16, 24, 32, 48, 64, 96, 128) & 0.740 & 0.935 & 0.795 & 0.994 & 0.718 & 0.930 \\
           
            \bottomrule
        \end{tabular}
        \begin{tablenotes}
            \footnotesize
             \item 特征计算的具体方法和尺度的定义详见前文所述。
        \end{tablenotes}
        \end{threeparttable}
       
\end{table*}

\subsection{训练策略的有效性分析}

针对类别不平衡问题，实验首先比较了两种数据重采样策略：一是仅对角点进行重采样，使其数量与边缘点一致；二是同时对角点和边缘点进行重采样，使其数量与面点一致。表~\ref{tab:additional_ablation_study} 的结果显示，两种策略在最终的检测精度上差异极小。然而，第二种策略会使训练数据集的规模呈指数级增长，导致训练时长大幅延长且内存开销巨大。基于效率优先的原则，本文最终采用了“仅重采样角点”的策略。

此外，受边缘检测领域中“平滑边缘（Smooth-Edge）”策略（即采样尖锐边缘附近的点作为平滑边缘点以提升检测鲁棒性）的启发，本实验尝试了类似的“平滑角点（Smooth-Corner）”策略，将角点邻域内的点也标记为角点参与训练。然而，表~\ref{tab:additional_ablation_study} 的实验结果表明，引入平滑角点后，模型的各项指标反而出现了下滑。深入分析认为，边缘特征通常具有一定的连续性，因此平滑操作有助于增强边缘的连通性；而角点在几何拓扑上属于“奇点”，具有高度的局部化特征。平滑操作人为地模糊了角点与边缘点之间的特征界限，混淆了分类边界，反而增加了分类器的学习难度。因此，本文在最终方案中摒弃了平滑角点策略，坚持对角点进行精确的单点监督。

\begin{table*}[ht]
    \centering
    \small
    \setlength{\tabcolsep}{10pt}
    \caption{关于训练策略的消融实验结果。}
    \label{tab:additional_ablation_study}
    \begin{tabular}{llcccccc}
            \toprule
            & & Precision($\uparrow$) & Recall($\uparrow$) & F1($\uparrow$)& Accuracy($\uparrow$)& IoU($\uparrow$)& MCC($\uparrow$)\\
            \midrule
   
            \multirow{2}{*}{重采样策略}
            & 边缘与角点 & 0.741& \textbf{0.941}& 0.795& \textbf{0.994}& 0.721& 0.923\\
            & 仅角点 & \textbf{0.743}& 0.940& \textbf{0.797}& \textbf{0.994}& \textbf{0.722}& \textbf{0.929}\\
           
            \midrule
           
            \multirow{2}{*}{平滑角点策略}
            & 使用平滑 & 0.706& 0.874& 0.736& 0.989& 0.670& 0.868\\
            & 不使用平滑 & \textbf{0.743}& \textbf{0.940}& \textbf{0.797}& \textbf{0.994}& \textbf{0.722}& \textbf{0.929}\\
           
            \bottomrule
        \end{tabular}
\end{table*}

\section{本章小结}

本章通过一系列严谨的对比实验与消融分析，系统评估了 NECDNet 在点云几何特征检测任务中的性能。实验结果表明，NECDNet 在边缘与角点检测任务上显著优于现有的传统分析方法及其他深度学习模型，展现出更强的特征提取能力和分类精度。通过对法向量精度的深入分析，揭示了高质量几何先验信息对尖锐特征检测的重要性。同时，消融实验验证了本文在特征维度选择、多尺度架构设计及数据平衡策略上的设计合理性，确立了最佳的模型配置。本章的实验结论充分证明了 NECDNet 能够有效地从复杂点云中提取出高质量的几何特征，为机器人抓取规划等下游任务提供了可靠的感知基础。